Amharic, a federal working language of Ethiopia with about 57 million first- and second-language speakers, remains severely underrepresented in natural language processing research. In this paper, we present a systematic study of sentiment analysis on Amharic Twitter data using the AfriSenti benchmark. We make three contributions: (1) a structured Ethiopic character normalization pipeline that unifies orthographically equivalent character families and improves the TF-IDF baseline's test macro F1 by 1.7 points in an ablation; (2) a comparison of a TF-IDF and Logistic Regression baseline against fine-tuned Afro-XLM-R, a multilingual transformer pre-trained on 17 African languages; and (3) a native-speaker qualitative error analysis of the baseline and the transformer, identifying five linguistically grounded failure patterns. Our fine-tuned model achieves a macro F1 of \textbf{0.5543 $\pm$ 0.0068} across three random seeds, a \textbf{+9.03 point improvement} over the TF-IDF baseline (0.4640). The error analysis separates three patterns on which the baseline fails and the transformer succeeds (journalistic framing with embedded sentiment words, Amharic negation, and culturally loaded religious language) from two on which both models fail (news-framed atrocity reporting and Amharic praise song structure); a feature analysis additionally shows the baseline learning a political name as a strong negative cue.
